% !TEX root = task-1-3-6.tex
% Editable task bank. Append new BankTask blocks at the end of the document.
\documentclass[12pt,a4paper]{article}
\usepackage{physics-textbook}
\usepackage{graphicx,needspace}
\usetikzlibrary{decorations.pathmorphing}
\newif\ifShowSolutions
\ifdefined\TutorEdition
  \ShowSolutionstrue
\else
  \ShowSolutionsfalse % Change to \ShowSolutionstrue to display solutions.
\fi
\newcounter{banktask}
\newsavebox{\BankTaskBox}
\newcommand{\BankTaskContent}[3]{%
  \noindent\textbf{\thebanktask.}\quad #2\par
  \ifShowSolutions
    \Needspace{50mm}% Keep the solution heading with its opening text or diagram.
    #3\par
    \ifdefined\TutorEdition
      \ifcsname TutorSource#1\endcsname
        \nopagebreak\smallskip{\small\csname TutorSource#1\endcsname\par}%
      \fi
    \fi
  \fi
}
\newcommand{\BankTask}[3]{%
  \refstepcounter{banktask}\par\medskip
  \sbox{\BankTaskBox}{\begin{minipage}{\linewidth}%
    \BankTaskContent{#1}{#2}{#3}\end{minipage}}%
  \ifdim\dimexpr\ht\BankTaskBox+\dp\BankTaskBox\relax>\dimexpr\textheight-20mm\relax
    \Needspace{35mm}\BankTaskContent{#1}{#2}{#3}%
  \else
    \Needspace{\dimexpr\ht\BankTaskBox+\dp\BankTaskBox\relax}%
    \noindent\usebox{\BankTaskBox}\par
  \fi
  \medskip
}
\begin{document}
\BookSubsection{1.3.6}{Бросок под углом к горизонту}


% Source sdamgia-4962: https://phys-ege.sdamgia.ru/problem?id=4962
\BankTask{sdamgia-4962}{
Тело, брошенное с горизонтальной поверхности со скоростью $v$ под углом $\alpha $ к горизонту, в течение $t$ секунд поднимается над горизонтом, а затем снижается и падает на расстоянии $S$ от точки броска. Сопротивление воздуха пренебрежимо мало.

Установите соответствие между физическими величинами и формулами, по которым их можно рассчитать. 

\TaskMatchingColumns{Физические величины}{А) время подъема $t$ на максимальную высоту \par 
Б) расстояние $S$ от точки броска до точки падения}{Формулы}{$1$) $\frac{v ^2 \sin  ^2 \alpha }{2g}$\par 
$2$) $\frac{v \cos  ^2 \alpha }{g}$\par 
$3$) $\frac{v ^2 \sin  2 \alpha }{g}$\par 
$4$) $\frac{v \sin  \alpha }{g}$}

\begin{center}\small\renewcommand{\arraystretch}{1.25}\begin{tabular}{|p{0.450\linewidth}|p{0.450\linewidth}|}\hline
$A$ & Б\\\hline
 & \\\hline
\end{tabular}\end{center}
}{
\smallskip\noindent\textbf{Решение}\par
В пренебрежении силой сопротивления воздуха, изменение вертикальной проекции тела со временем описывается выражением:

\[v _y ( t ) = v \sin  \alpha  - gt.\]

Окончанию подъема соответствует момент времени, когда вертикальная проекция скорости обращается в ноль, то есть:

\[v _y ( t ) =0 \Longleftrightarrow  t= \frac{v \sin  \alpha }{g} .\]

Время подъема равно времени опускания, поэтому время полета равно: 

\[T=2t= \frac{2 v \sin  \alpha }{g} .\]

Все это время тело по горизонтали двигается с постоянной скоростью $v \cos  \alpha  .$ Следовательно, расстояние $S$ от точки броска до точки падения равно:

\[S= v \cos  \alpha  T = \frac{v ^2 \sin  2 \alpha }{g} .\]

Ответ: $43$.
}

% Source sdamgia-5172: https://phys-ege.sdamgia.ru/problem?id=5172
\BankTask{sdamgia-5172}{
Тело, брошенное с горизонтальной поверхности со скоростью $v$ под углом $\alpha $ к горизонту, в течение времени $t$ поднимается на максимальную высоту h над горизонтом. Сопротивление воздуха пренебрежимо мало.

Установите соответствие между физическими величинами и формулами, по которым их можно рассчитать. 

\TaskMatchingColumns{Физические величины}{А) время подъема $t$ на максимальную высоту \par 
Б) максимальная высота h над горизонтом}{Формулы}{$1$) $\frac{v ^2 \sin  ^2 \alpha }{2g}$\par 
$2$) $\frac{v \cos  ^2 \alpha }{g}$\par 
$3$) $\frac{v ^2 \sin  2 \alpha }{2g}$\par 
$4$) $\frac{v \sin  \alpha }{g}$}

\begin{center}\small\renewcommand{\arraystretch}{1.25}\begin{tabular}{|p{0.450\linewidth}|p{0.450\linewidth}|}\hline
$A$ & Б\\\hline
 & \\\hline
\end{tabular}\end{center}
}{
\smallskip\noindent\textbf{Решение}\par
В пренебрежении силой сопротивления воздуха, изменение вертикальной проекции тела со временем описывается выражением:

\[v _y ( t ) = v \sin  \alpha  - gt.\]

Окончанию подъема соответствует момент времени, когда вертикальная проекция скорости обращается в ноль, то есть:

\[v _y ( t ) =0 \Rightarrow t= \frac{v \sin  \alpha }{g} .\]

Вертикальная координата тела зависит от времени согласно

\[y ( t ) = v \sin  \alpha  t - \frac{gt ^2}{2} .\]

Следовательно, к моменту окончания подъема, тело достигнет высоты:

\[h= v \sin  \alpha  \frac{v \sin  \alpha }{g} - \frac{g ( v \sin  \alpha  /g ) ^2}{2} = \frac{v ^2 \sin  ^2 \alpha }{2g} .\]

Ответ: $41$.
}

% Source sdamgia-5207: https://phys-ege.sdamgia.ru/problem?id=5207
\BankTask{sdamgia-5207}{
Тело, брошенное с горизонтальной поверхности со скоростью $v$ под углом $\alpha $ к горизонту, через некоторое время $t$ падает на расстоянии $S$ от точки броска. Сопротивление воздуха пренебрежимо мало.

Установите соответствие между физическими величинами и формулами, по которым их можно рассчитать. 

\TaskMatchingColumns{Физические величины}{А) время полета $t$\par 
Б) расстояние $S$ от точки броска до точки падения}{Формулы}{$1$) $\frac{v ^2 \sin  ^2 \alpha }{2g}$\par 
$2$) $\frac{2 v \sin  \alpha }{g}$\par 
$3$) $\frac{v \sin  \alpha }{2g}$\par 
$4$) $\frac{v ^2 \sin  2 \alpha }{g}$}

\begin{center}\small\renewcommand{\arraystretch}{1.25}\begin{tabular}{|p{0.450\linewidth}|p{0.450\linewidth}|}\hline
$A$ & Б\\\hline
 & \\\hline
\end{tabular}\end{center}
}{
\smallskip\noindent\textbf{Решение}\par
В пренебрежении силой сопротивления воздуха, изменение вертикальной проекции тела со временем описывается выражением:

\[v _y ( t ) = v \sin  \alpha  - gt.\]

Окончанию подъема соответствует момент времени, когда вертикальная проекция скорости обращается в ноль, то есть:

\[v _y ( \tau ) = \Longleftrightarrow  \tau= \frac{v \sin  \alpha }{g} .\]

Время подъема равно времени опускания, поэтому время полета равно: 

\[t=2\tau= \frac{2 v \sin  \alpha }{g} .\]

Все это время тело по горизонтали двигается с постоянной скоростью $v \cos  \alpha  .$ Следовательно, расстояние $S$ от точки броска до точки падения равно:

\[S= v \cos  \alpha  T = \frac{v ^2 \sin  2 \alpha }{g} .\]

Ответ: $24$.
}

% Source sdamgia-5312: https://phys-ege.sdamgia.ru/problem?id=5312
\BankTask{sdamgia-5312}{
Тело, брошенное с горизонтальной поверхности со скоростью $v$ под углом $\alpha $ к горизонту, поднимается над горизонтом на максимальную высоту h, а затем падает на расстоянии $S$ от точки броска. Сопротивление воздуха пренебрежимо мало.

Установите соответствие между физическими величинами и формулами, по которым их можно рассчитать. 

\TaskMatchingColumns{Физические величины}{А) Максимальная высота h над горизонтом \par 
Б) Расстояние $S$ от точки броска до точки падения}{Формулы}{$1$) $\frac{v ^2 \sin  ^2 \alpha }{2g}$\par 
$2$) $\frac{v ^2 \cos  ^2 \alpha }{g}$\par 
$3$) $\frac{v ^2 \sin  2 \alpha }{g}$\par 
$4$) $\frac{v ^2 \sin  \alpha }{g}$}

\begin{center}\small\renewcommand{\arraystretch}{1.25}\begin{tabular}{|p{0.450\linewidth}|p{0.450\linewidth}|}\hline
$A$ & Б\\\hline
 & \\\hline
\end{tabular}\end{center}
}{
\smallskip\noindent\textbf{Решение}\par
В пренебрежении силой сопротивления воздуха, изменение вертикальной проекции тела со временем описывается выражением:

\[v _y ( t ) = v \sin  \alpha  - gt.\]

Окончанию подъема соответствует момент времени, когда вертикальная проекция скорости обращается в ноль, то есть:

\[v _y ( t ) =0 \Longleftrightarrow  t= \frac{v \sin  \alpha }{g} .\]

Вертикальная координата тела зависит от времени согласно:

\[y ( t ) = v \sin  \alpha  t - \frac{gt ^2}{2} .\]

Следовательно, к моменту окончания подъема тело достигнет высоты:

\[h= v \sin  \alpha  \frac{v \sin  \alpha }{g} - \frac{g ( v \sin  \alpha  /g ) ^2}{2} = \frac{v ^2 \sin  ^2 \alpha }{2g} .\]

Время подъема равно времени опускания, поэтому время полета равно: 

\[T=2t= \frac{2 v \sin  \alpha }{g} .\]

Все это время тело по горизонтали двигается с постоянной скоростью $v \cos  \alpha  .$ Следовательно, расстояние $S$ от точки броска до точки падения равно:

\[S= v \cos  \alpha  T = \frac{v ^2 \sin  2 \alpha }{g} .\]

Ответ: $13$.
}

% Source sdamgia-7137: https://phys-ege.sdamgia.ru/problem?id=7137
\BankTask{sdamgia-7137}{
Камень брошен вверх под углом к горизонту. Сопротивление воздуха пренебрежимо малó. Как меняются с набором высоты потенциальная энергия камня в поле тяжести и ускорение камня?

Для каждой величины определите соответствующий характер изменения.

$1$. Увеличивается.

$2$. Уменьшается.

$3$. Не изменяется.

 

\begin{center}\small\renewcommand{\arraystretch}{1.25}\begin{tabular}{|p{0.450\linewidth}|p{0.450\linewidth}|}\hline
Потенциальная энергия

камня & Ускорение камня\\\hline
 & \\\hline
\end{tabular}\end{center}
}{
\smallskip\noindent\textbf{Решение}\par
Потенциальная энергия камня увеличивается, потому что Е = mgh. Ускорение камня не меняется и равно $g$.

Ответ: $13$.
}

% Source sdamgia-50759: https://phys-ege.sdamgia.ru/problem?id=50759
\BankTask{sdamgia-50759}{
Камень брошен вверх под углом к горизонту. Сопротивление воздуха пренебрежимо мало. Как меняются по мере подъёма камня модуль ускорения камня и его кинетическая энергия? Для каждой величины определите соответствующий характер её изменения:

$1$. Увеличится

$2$. Уменьшится

$3$. Не изменится

 

\begin{center}\small\renewcommand{\arraystretch}{1.25}\begin{tabular}{|p{0.450\linewidth}|p{0.450\linewidth}|}\hline
Модуль ускорения камня & Кинетическая энергия камня\\\hline
 & \\\hline
\end{tabular}\end{center}
}{
\smallskip\noindent\textbf{Решение}\par
При отсутствии сопротивления воздуха камень движется под действием силы тяжести с ускорением свободного падения. Его модуль не меняется при движении тела ($3$).

По мере подъема камня модуль его скорости уменьшается, поэтому и кинетическая энергия уменьшается ($2$).

Ответ: $32$.
}

% Source sdamgia-8874: https://phys-ege.sdamgia.ru/problem?id=8874
\BankTask{sdamgia-8874}{
Точечное тело брошено под углом $45^\circ$ к горизонту со скоростью $20\,\text{м/с}$. Пренебрегая сопротивлением воздуха, определите модуль скорости этого тела через $0{,}47\,\text{с}$ после броска. Ответ выразите в метрах в секунду округлите до целого числа.
}{
\smallskip\noindent\textbf{Решение}\par
Запишем проекции скоростей на горизонтальную и вертикальную оси: $v _x= v _0 \cdot  \cos  45 ^\circ. v _y= v _0 \cdot  \sin  45 ^\circ - gt.$

Через $0{,}47\,\text{с}$ скорость тела будет равна:

\[v = \sqrt{v _x ^2 + v _y ^2} = \sqrt{v _0 ^2 \cos  ^2 45 ^\circ + v _0 ^2 \sin  ^2 45 ^\circ - 2 v _0gt \sin  45 ^\circ + ( gt ) ^2} =\]

\[= \sqrt{v _0 ^2 ( \cos  ^2 45 ^\circ + \sin  ^2 45 ^\circ ) - 2 v _0gt \sin  45 ^\circ + ( gt ) ^2}\]

\[= \sqrt{v _0 ^2 - 2 v _0gt \sin  45 ^\circ + ( gt ) ^2} = \sqrt{400 - 2 \cdot  20 \cdot  10 \cdot  0{,}47 \cdot  \frac{\sqrt{2}}{2} + 100 \cdot  0{,}47 \cdot  0{,}47} \approx17\,\text{м/с}.\]

Ответ: $17$.
}

% Source sdamgia-9517: https://phys-ege.sdamgia.ru/problem?id=9517
\BankTask{sdamgia-9517}{
Камень массой $40\,\text{г}$ брошен под углом $60^\circ$ к горизонту. Начальная кинетическая энергия камня равна $2\,\text{Дж}$. Чему равен модуль импульса камня в верхней точке траектории его движения? Сопротивление воздуха пренебрежимо мало.
}{
\smallskip\noindent\textbf{Решение}\par
При отсутствии сопротивления воздуха камень движется только под действием силы тяжести с ускорением свободного падения. Его проекция на ось Ox равна $g_x=0,$ следовательно, движение камня в горизонтальном направлении равномерное $v _x= v _{0x}.$ Проекции скорости в верхней точке траектории равны $v _{0x}= v _0 \cos  \alpha  ,$ $v _y=0,$ так как в верхней точке траектории скорость камня направлена горизонтально. Таким образом, скорость камня в верхней точке траектории равна $v = v _x= v _0 \cos  \alpha  .$

Кинетическая энергия камня в момент броска $E_{\text{к}}= \frac{m v _0 ^2}{2} ,$ откуда его начальная скорость $v _0= \sqrt{\frac{2E_{\text{к}}}{m}} .$

Тогда импульс камня в верхней точке траектории равен:

\[p=m v =m v _0 \cos  \alpha  = \sqrt{2E_{\text{к}}m} \cos  \alpha  = \sqrt{2 \cdot  2 \cdot  0{,}04} \cdot  \cos  60 ^\circ =0{,}2\,\text{кг} \cdot  \text{м/с} .\]

Ответ: $0{,}2\,\text{кг} \cdot  \text{м/с} .$
}

% Source sdamgia-37878: https://phys-ege.sdamgia.ru/problem?id=37878
\BankTask{sdamgia-37878}{
Маленький камень бросили с поверхности Земли под углом $45^\circ$ к горизонту, он упал обратно на землю на расстоянии $20$ метров от точки броска. Найти максимальную высоту подъема камня при полете. Сопротивлением воздуха пренебречь.
}{
\smallskip\noindent\textbf{Решение}\par
Камень движется под действием силы тяжести с ускорением свободного падения, направленным вертикально вниз. Проекция ускорения на ось Ox равна нулю, поэтому движение по этой оси равномерное; на ось Oy равна –$g$, поэтому движение по этой оси равноускоренное.

При равноускоренном движении проекция перемещения: 

\[s_y = v _{0y}t - \frac{gt ^2}{2} =0.\]

Откуда время полета тела равно: 

\[t_{\text{пол}} = \frac{2 v _{0y}}{g} = \frac{2 v _0 \sin  \alpha }{g} .\]

При равномерном движении по оси Ox дальность полета $L= v _{0x}t_{\text{пол}}= v _0t_{\text{пол}} \cos  \alpha  ,$ откуда начальная скорость при броске равна: 

\[v _0= \sqrt{\frac{Lg}{2 \sin  \alpha  \cos  \alpha }} .\]

При подъеме камня на максимальную высоту проекция $v _y=0.$ При равноускоренном движении высота подъема равна: 

\[h_{max}= \frac{v _{0y} ^2}{2g} .\]

Тогда получаем:

\[h_{max}= \frac{v _0 ^2 \sin  ^2 \alpha }{2g} = \frac{L \tan  \alpha }{4} = \frac{20}{4} =5\,\text{м} .\]

Ответ: $5\,\text{м}$.
}

% Source sdamgia-38204: https://phys-ege.sdamgia.ru/problem?id=38204
\BankTask{sdamgia-38204}{
Небольшой камень, брошенный с ровной горизонтальной поверхности Земли под углом $60^\circ$ к горизонту, достиг максимальной высоты, равной $5\,\text{м}$. Сколько времени прошло от момента броска до того момента, когда скорость камня стала горизонтальной? Сопротивлением воздуха пренебречь.
}{
\smallskip\noindent\textbf{Решение}\par
При движении тела, брошенного под углом к горизонту, на него действует только сила тяжести (по условию сопротивлением воздуха пренебречь). Данная сила является причиной ускорения свободного падения, направленного вертикально вниз.

Выберем ось Oy направленную вертикально вверх. Проекция начальной скорости камня на данную ось $v _{0y}= v _0 \sin  \alpha  ,$ проекция ускорения $g_y= - g.$ Движение по оси Oy равноускоренное. Тогда проекция скорости на эту ось в любой точке траектории равна $v _y= v _{0y} - gt,$ а проекция перемещения:

\[s_y= \frac{v _y ^2 - v _{0y} ^2}{- 2g} .\]

В точке максимального подъема камня его скорость направлена горизонтально, следовательно, $v _y=0.$

Применим формулы для данных условий: 

$0= v _0 \sin  \alpha  - gt$ и $H= \frac{v _0 ^2 \sin  ^2 \alpha }{2g} .$

При объединении этих формул, получаем:

\[v _0 \sin  \alpha  =gt;\]

\[H= \frac{gt ^2}{2} .\]

Тогда время подъема камня $t= \sqrt{\frac{2H}{g}} .$

При подстановке значений получаем: $t= \sqrt{\frac{2 \cdot  5}{10}} =1\,\text{с}.$

Ответ: $1\,\text{с}$.
}


% Source F4CD91: https://ege.fipi.ru/bank/questions.php?search=1&proj=BA1F39653304A5B041B656915DC36B38&qid=F4CD91
\BankTask{F4CD91}{
Тело, брошенное с горизонтальной поверхности Земли со скоростью $v$ под углом $\alpha$ к горизонту, поднимается на максимальную высоту $h$ над земной поверхностью и затем падает на землю. Сопротивлением воздуха пренебречь.
\TaskMatchingColumns[.53\linewidth]{Физические величины}{
А) максимальная высота $h$ над горизонтом\par
Б) модуль вертикальной проекции скорости тела непосредственно перед падением на землю
}{Формулы}{
1) $\frac{v^2\sin^2\alpha}{2g}$\par
2) $v\sin\alpha$\par
3) $\frac{v^2\sin2\alpha}{2g}$\par
4) $v\cos\alpha$\par
}\begin{center}\begin{tabular}{|c|c|}\hline А & Б \\\hline\rule{0pt}{1.2em}\hspace{1cm} & \hspace{1cm} \\\hline\end{tabular}\end{center}
}{
\smallskip\noindent\textbf{Решение от GPT}\par
Начальная вертикальная скорость $v_{0y}=v\sin\alpha$. Из $v_y^2=v_{0y}^2-2gh$ при $v_y=0$ получаем $h=v^2\sin^2\alpha/(2g)$ (формула $1$). При возвращении на исходный уровень $v_y=-v\sin\alpha$, поэтому $|v_y|=v\sin\alpha$ (формула $2$). Ответ: $12$.
}

% Source 32F41D: https://ege.fipi.ru/bank/questions.php?search=1&proj=BA1F39653304A5B041B656915DC36B38&qid=32F41D
\BankTask{32F41D}{
Шарик, брошенный с поверхности Земли под углом $20\,{}^\circ$ к горизонту
с начальной скоростью ${\vec{\upsilon}}_0$, поднялся на максимальную высоту $H$ и пролетел
в горизонтальном направлении расстояние $L$. Что произойдёт с дальностью полёта шарика и временем его полёта, если шарик бросить с той же
по величине начальной скоростью под углом $30\,{}^\circ$ к горизонту?
Сопротивлением воздуха пренебречь.
\begin{enumerate}
\item увеличится
\item уменьшится
\item не изменится
\end{enumerate}
\TaskChangeTable{Дальность полёта шарика}{Время полёта шарика}
}{
\smallskip\noindent\textbf{Решение от GPT}\par
$L=v_0^2\sin2\alpha/g$, $t=2v_0\sin\alpha/g$. При росте угла от $20^\circ$ до $30^\circ$ возрастают и $\sin2\alpha$, и $\sin\alpha$, следовательно, обе величины увеличиваются. Ответ: $11$.
}

% Source A5cD6A: https://ege.fipi.ru/bank/questions.php?search=1&proj=BA1F39653304A5B041B656915DC36B38&qid=A5cD6A
\BankTask{A5cD6A}{
Шарик, брошенный с поверхности Земли под углом $\alpha$ к горизонту с начальной скоростью ${\vec{\upsilon}}_0$, поднялся на максимальную высоту $H$ и пролетел в горизонтальном направлении расстояние $L$. Что произойдёт с максимальной высотой, на которую поднимется шарик, и его ускорением, если уменьшить начальную скорость шарика, не меняя $\alpha$? Сопротивлением воздуха пренебречь.
\begin{enumerate}
\item увеличится
\item уменьшится
\item не изменится
\end{enumerate}
\TaskChangeTable{Максимальная высота подъёма}{Ускорение шарика}
}{
\smallskip\noindent\textbf{Решение от GPT}\par
$H=v_0^2\sin^2\alpha/(2g)$ уменьшается при уменьшении $v_0$ с неизменным углом. Ускорение $\vec a=\vec g$ не меняется. Ответ: $23$.
}

% Source 25383: https://phys-ege.sdamgia.ru/problem?id=25383
\BankTask{25383}{
Человек ростом h  =  $1{,}6\,\text{м}$, стоя на земле, кидает мяч из-за головы и хочет перебросить его через забор высотой H  =  $4{,}8\,\text{м}$, находящийся на расстоянии l  =  $6{,}4\,\text{м}$ от него. Определите модуль скорости, с которой необходимо бросить мяч, чтобы он перелетел через забор, коснувшись его в верхней точке своей траектории? Сопротивлением воздуха пренебречь.
}{
\smallskip\noindent\textbf{Решение}\par
Известно, что в верхней точке траектории тела, брошенного под углом к горизонту, скорость направлена горизонтально. Без действия сопротивления воздуха тело движется с ускорением свободного падения. Поэтому по оси Ox движение является равномерным (горизонтальная проекция скорости одинакова на всем пути), по оси Oy движение равноускоренное, при этом в верхней точки траектории вертикальная проекция скорости равна $0$.
Применим закон сохранения энергии для движения мячика от точки броска до верхней точки траектории

\[mgh + \frac{m v _0 ^2}{2} =mgH + \frac{m v _x ^2}{2} .\]

Учитывая, что $v _0 ^2 = v_{0x} ^2 + v_{0y} ^2 ,$ найдем вертикальную проекцию начальной скорости:

\[v_{0y}= \sqrt{2g ( H - h )} = \sqrt{2 \cdot 10 \cdot ( 4{,}8 - 1{,}6 )} =8 \,\text{м/с}.\]

Скорость тела при равноускоренном движении по вертикали равна $v _y= v_{0y} - gt,$ откуда время подъема до верхней точки траектории $t= \frac{v_{0y}}{g} = \frac{8}{10} =0{,}8 \,\text{с}.$
Зная расстояние, которое мяч пролетел по горизонтали, найдем горизонтальную составляющую начальной скорости $v _x= \frac{l}{t} = \frac{6{,}4}{0{,}8} =8 \,\text{м/с}.$
Тогда начальная скорость броска мяча

\[v _0= \sqrt{v_{0x}^2+v_{0y}^2} =8 \sqrt{2} \approx 11{,}3\,\text{м/с}.\]

Ответ: $11{,}3\,\text{м/с}$.
}

% Source 39272: https://phys-ege.sdamgia.ru/problem?id=39272
\BankTask{39272}{
Пластилиновый шарик в момент t  =  $0$ бросают с горизонтальной поверхности Земли под углом $\alpha$ к горизонту. Одновременно с некоторой высоты над поверхностью Земли начинает падать из состояния покоя другой такой же шарик. Шарики абсолютно неупруго сталкиваются в воздухе. Сразу после столкновения скорость шариков направлена горизонтально. Время от столкновения шариков до их падения на Землю равно $\tau$. С какой начальной скоростью $\upsilon_{0}$ был брошен первый шарик? Сопротивлением воздуха пренебречь.
}{
\smallskip\noindent\textbf{Решение}\par
\begin{center}\begin{tikzpicture}[task graph,>=stealth]\draw[->](0,0)--(8.8,0)node[right]{$x$};\draw[->](0,0)--(0,4.8)node[above]{$y$};\node[below left]at(0,0){$0$};\draw[domain=0:8.4,samples=60]plot(\x,{1.333333*\x-.15873016*\x*\x});\draw[->](0,0)--(.8,1.066667)node[above left]{$\vec v_0$};\draw(.6,0)arc(0:53.1301:.6)node[midway,right]{$\alpha$};\coordinate(B)at(2.1,2.1);\coordinate(A)at(2.1,3.4);\fill(B)circle(1.7pt)node[left]{$B$};\fill(A)circle(1.7pt)node[left]{$A$};\draw(A)--(B);\draw[->,thick](B)--(3.8,3.233333)node[above]{$m\vec v_1$};\draw[->,thick](B)--(3.8,2.1)node[right]{$2m\vec u_0$};\draw[->,thick](B)--(2.1,.966667)node[right]{$m\vec v_2$};\draw[dashed](3.8,3.233333)--(3.8,2.1);\draw[dashed](2.1,.966667)--(3.8,2.1);\end{tikzpicture}\end{center}Обоснование.
$1$.  Выберем инерциальную систему отсчёта, связанную с Землёй. За начало отсчёта координат примем первоначальное положение первого шарика.
$2$.  Шарики будем считать материальными точками.
$3$.  Так как сопротивлением воздуха можно пренебречь, то движение шариков можно считать свободным падением.
$4$.  Считаем время взаимодействия шариков при неупругом столкновении малым. Следовательно, импульсом внешней силы (силы тяжести) за это время можно пренебречь. Значит, импульс системы двух шариков при столкновении сохраняется. 

Перейдем к решению.
$1$.  Первый шарик начинает движение из начала координат, а второй  — из точки $A$. До и после столкновения, происходящего в точке $B$, шарики свободно падают. Следовательно, до столкновения для первого шарика

$y_1 ( t ) = v _0y t - \frac{gt ^2}{2} = v _0 \sin \alpha \cdot t - \frac{gt ^2}{2} ,$ 
$v _1y ( t ) = v _0 \sin \alpha - gt,$

а для второго шарика $v _2y ( t ) = - gt.$
$2$.  Шарики сталкиваются в момент $t_{1}$, при этом импульс системы двух шариков сохраняется: $m \vec v _1 + m \vec v _2 = 2m \vec u_0,$ а скорость $\vec u_0$ шариков после удара согласно условию горизонтальна. Следовательно,

\[v _1y ( t_1 ) + v _2y ( t_1 ) = 0 \Longleftrightarrow ( v _0 \sin \alpha - gt_1 ) + ( - gt_1 ) = 0 \Longleftrightarrow t_1 = \frac{v _0 \sin \alpha}{2g} .\]

$3$.  Столкновение шариков происходит на высоте

\[h = y_1 ( t_1 ) = v _0 \sin \alpha \cdot t_1 - \frac{gt_1 ^2}{2} = \frac{v _0 ^2 \sin ^2 \alpha}{2g} - \frac{v _0 ^2 \sin ^2 \alpha}{8g} = \frac{3}{8} \cdot \frac{v _0 ^2 \sin ^2 \alpha}{g} .\]

$4$.  Скорость $\vec u_0$ шариков после удара горизонтальна, поэтому интервал времени $\tau$ от столкновения шариков до их падения на Землю находится из условия $h = \frac{g \tau ^2}{2} ,$ откуда $\tau = \sqrt{\frac{2h}{g}} = \sqrt{3} \cdot \frac{v _0 \sin \alpha}{2g} .$ Таким образом, $v _0 = \frac{2g \tau}{\sqrt{3} \sin \alpha} .$

Ответ: $v _0 = \frac{2g \tau}{\sqrt{3} \sin \alpha} .$
}

% Source 38860: https://phys-ege.sdamgia.ru/problem?id=38860
\BankTask{38860}{
\begin{center}\begin{tikzpicture}[task graph,>=stealth]\draw(-.82,0)--(4,0);\draw(-.6,0)--(-.6,3.3);\draw[<->](-.6,0)--(-.6,3.3)node[midway,left]{\(H\)};\draw(-.6,3.3)--(0,3.3);\draw[<->](0,0)--(0,1.96)node[midway,left]{\(h\)};\draw[dashed](0,1.96)--(0,3.3);\draw[fill=gray!15](0,3.3)circle(.21);\draw(-.46,2.23)--(.46,1.69);\draw(0,2.33)arc(90:149:.37)node[midway,above left]{\(\alpha\)};\draw[dashed](0,1.96) parabola bend(2.35,3.02)(3.7,2.72);\draw[<->](2.35,0)--(2.35,3.02)node[midway,left]{\(h_1\)};\end{tikzpicture}\end{center}Шарик падает из состояния покоя и на высоте h  =  $1\,\text{м}$ от поверхности Земли абсолютно упруго ударяется о доску, расположенную под углом $\alpha =30 ^\circ$ к горизонту (см. рисунок). Модуль скорости шарика после отскока от доски равен модулю скорости шарика непосредственно перед ударом. После удара шарик поднялся на максимальную высоту h$_{1}$  =  $1{,}25\,\text{м}$ от поверхности Земли. Определите с какой высоты H над поверхностью Земли упал шарик, если после удара его скорость направлена под углом $2 \alpha$ к вертикали. Сопротивлением воздуха пренебречь. 
Обоснуйте применимость используемых законов к решению задачи.
}{
\smallskip\noindent\textbf{Решение}\par
Обоснование
Систему отсчета, связанную с Землей, будем считать инерциальной. Шарик движется поступательно, поэтому его можно принять за материальную точку. По условию задачи силой сопротивления воздуха можно пренебречь. Поэтому шарик движется в поле Земного тяготения. Следовательно, для описания движения тел до столкновения с доской можно использовать закон сохранения энергии
Поскольку при ударе о доску модуль шарика не изменился, то можно считать, что удар абсолютно упругий, следовательно, угол отскока равен углу при падении и равен $\alpha .$
Движение шарика после отскока от доски происходит под действием только силы тяжести, поэтому он движется с ускорением свободного падения.

Перейдем к решению. \begin{center}\begin{tikzpicture}[task graph,>=stealth]\draw(-.82,0)--(1.3,0);\draw(-.6,0)--(-.6,3.3);\draw[<->](-.6,0)--(-.6,3.3)node[midway,left]{\(H\)};\draw(-.6,3.3)--(0,3.3);\draw[<->](0,0)--(0,1.96)node[midway,left]{\(h\)};\draw[dashed](0,1.96)--(0,3.3);\draw[fill=gray!15](0,3.3)circle(.21);\draw(-.46,2.23)--(.46,1.69);\draw(0,2.33)arc(90:149:.37)node[midway,above left]{\(\alpha\)};\draw[->](0,1.96)--(.57,2.84)node[above]{\(\vec n\)};\draw[->](0,1.96)--(1.28,2.54);\draw[dashed](0,1.96)--(1.36,1.96);\draw(1.02,1.96)arc(0:24.4:1.02)node[midway,right]{\(\beta\)};\draw(0,2.54)arc(90:57:.58)node[midway,above right]{\(\alpha\)};\end{tikzpicture}\end{center}
Рассмотрим движение шарика до соударения о доску. В верхней точке он обладал потенциальной энергией $E_{\,\text{п}1}= mgH,$ в момент удара  — потенциальной $E_{\,\text{п}2}= mgh$ и кинетической $E_{\,\text{к}2}= \frac{m v ^2}{2} ,$ где m  — масса шарика, \ensuremath{\upsilon}  — скорость шарика в момент удара о доску. По закону сохранения энергии $E_{\,\text{п}1}=E_{\,\text{п}2} + E_{\,\text{к}2}.$ Тогда получаем:

$mgH=mgh + \frac{m v ^2}{2} ,$
$H=h + \frac{v ^2}{2g} . \qquad ( 1 )$

Из рисунка следует, что угол между скоростью при отскоке и горизонтом равен $\beta =90 ^\circ - 2 \alpha .$
Дальнейшее движение шарика происходило с ускорением свободного падения g. Проекция начальной скорости на вертикальную ось равна $v _y= v \sin \beta = v \sin ( 90 ^\circ - 2 \alpha ) = v \cos 2 \alpha .$ В точке максимального подъема $v_{y1}=0.$ Если выбрать в качестве начального уровня высоты точку отскока, то максимальная высота подъема при равноускоренном движении

\[h_1 - h= \frac{v_{y1} ^2 - v _y ^2}{- 2g} = \frac{v ^2 \cos ^2 2 \alpha}{2g} . \qquad ( 2 )\]

Объединяя формулы ($1$) и ($2$), получаем:

\[H=h + \frac{h_1 - h}{\cos ^2 2 \alpha} =1 + \frac{1{,}25 - 1}{0{,}25} =2\,\text{м}.\]

Ответ: $2\,\text{м}$.
}

% Source 25969: https://phys-ege.sdamgia.ru/problem?id=25969
\BankTask{25969}{
Пружинное ружье наклонено под углом $\alpha =30 ^\circ$ к горизонту. Энергия сжатой пружины равна $0{,}41\,\text{Дж}$. При выстреле шарик массой $m=50\,\text{г}$ проходит по стволу ружья расстояние b, вылетает и падает на расстоянии $L=1\,\text{м}$ от дула ружья в точку M, находящуюся с ним на одной высоте с дулом. (см. рис.). Найдите расстояние b. Трением в стволе и сопротивлением воздуха пренебречь.
Какие законы Вы использовали для описания движения шарика? Обоснуйте их применение к данному случаю.
\begin{center}\begin{tikzpicture}[task graph,>=stealth]\draw(0,0)--(4.6,0);\draw(0,0)--(4.6,2.75)--(10.2,2.75);\draw(0,0)--(-.25,.41)--(4.35,3.16)--(4.6,2.75);\draw[<->](1.2,1.2)--(4.31,3.06)node[midway,above]{\(b\)};\draw[<->](4.6,2.01)--(10.2,2.01)node[midway,below]{\(L\)};\draw(4.6,2.01)--(4.6,2.75)(10.2,2.01)--(10.2,2.75);\foreach \x in {4.7,4.85,...,10.1}{\draw(\x,2.75)--++(-.12,-.12);}\fill(10.2,2.75)circle(1.5pt)node[above right]{\(M\)};\draw(.72,0)arc(0:31:.72)node[midway,right]{\(\alpha\)};\begin{scope}[rotate=31]\draw plot[domain=.1:1.42,samples=120](\x,{.22+.13*sin(2727*\x)});\draw[fill=gray!15](1.63,.22)circle(.21)node[below right,xshift=2pt,yshift=-6pt]{\(m\)};\end{scope}\end{tikzpicture}\end{center}
}{
\smallskip\noindent\textbf{Решение}\par
Обоснование 
Систему отсчета, связанную с Землей, будем считать инерциальной. При движении пули по дулу пистолета на ее действуют силы тяжести, сила упругости и сила реакции опоры, не действуют сила трения (так как дуло можно считать абсолютно гладким) и сила сопротивления воздуха. В ИСО изменение полной механической энергии равно работе всех не потенциальных сил, действующих на тело. В данной задаче такой силой является сила реакции опоры $\vec N.$ Так как в процессе движения по стволу ружья сила реакции опоры $\vec N$ перпендикулярна поверхности, то ее работа равна нулю и $E_{\,\text{мех}}$ сохраняется в ИСО. 
При движении после вылета из трубки мы можем пренебречь действием силы сопротивления воздуха, значит, движение мяча происходит в поле Земного тяготения. Ось 0x направим горизонтально, ось 0y направим вертикально вверх. Тогда проекция вектора ускорения на ось 0x равна $0$, поэтому для описания движения по горизонтали можно использовать законы прямолинейного равномерного движения. Проекция вектора ускорения на ось 0y равна  — g, поэтому для описания движения по вертикали можно использовать законы прямолинейного равноускоренного движения.

Перейдем к решению.
$1$.  Рассмотрим движение шарика внутри трубки. По закону сохранения энергии $E_{\,\text{п}\,\text{пр}}=E_{\,\text{п}\,\text{ш}} + E_{\,\text{к}},$ где $E_{\,\text{п}\,\text{ш}}=mgh,$ $E_{\,\text{к}}= \frac{m v ^2}{2} .$ Откуда получаем:

\[E_{\,\text{пр}}=mgh + \frac{m v ^2}{2} .\]

По рисунку $h=b \cdot \sin \alpha .$
$2$.  Шарик вылетает из трубки со скоростью $v$ под углом $30 ^\circ$ к горизонту. При отсутствии сопротивления воздуха по горизонтальной оси движение равномерное, так как $g_x=0,$ по вертикальной оси, направленной вертикально вверх,  — равноускоренное с ускорением свободного падения, так как $g_y= - g.$ Такое движение описываем уравнениями $s_x= v _xt$ и $s_y= v _yt - \frac{gt ^2}{2} .$ Проекции начальной скорости при вылете из трубки $v _x= v \cos \alpha ,$ $v _y= v \sin \alpha .$ Учитывая, что в момент падения на горизонтальную плоскость $s_y=0,$ находим из приведенных формул время полета шарика: 

\[t= \frac{2 v \sin \alpha}{g} .\]

Тогда дальность полета шарика равна $L= \frac{v ^2 \sin 2 \alpha}{g} .$
$3$.  Объединяя записанные формулы, выразим искомую величину:

\[b= \frac{2E_{\,\text{пр}} \sin 2 \alpha - mgL}{2mg \sin \alpha \sin 2 \alpha} = \frac{2 \cdot 0{,}41 \cdot \sin 60 ^\circ - 0{,}05 \cdot 10 \cdot 1}{2 \cdot 0{,}05 \cdot 10 \cdot 0{,}5 \cdot \sin 60 ^\circ} \approx 0{,}485\,\text{м}=48{,}5\,\text{см}.\]

Ответ: $48{,}5\,\text{см}$.
}

\end{document}
